\vspace{-1mm}\section{Conclusion}
\label{sec:conclusion}
By bridging the gap between heuristic methods and the rigorous optimization framework of the Schr\"{o}dinger Bridge, this paper provides the first principled derivation of the log-density-ratio steering objective for activation steering. Experiments across four models and three alignment axes show that \Cobras consistently outperforms baselines in-distribution while preserving OOD performance (commonly seen in fixed-direction steering methods), thanks to its query-adaptive nature. \textcolor{blue}{The theoretical reduction is limited to a shared fixed bandwidth and does not establish optimality of the adaptive steering field at finite transport scale.} Limitations include the $\mathcal{O}(N_+ \cdot N_-)$ cost of pairwise geodesic distances and Sinkhorn, which may be prohibitive for very large contrastive sets, and the reliance on a fixed spherical geometry that may not hold for all architectures or layers. Promising future directions include hierarchical transport solvers to reduce computational overhead and extending the framework to learned or more general Riemannian manifolds.

